\begin{lemma}[Stored colors are invariant on old slots]
Let \(A^+\) be obtained by appending arbitrary final values to Tashkinov-series
data \(A\) as in \(\noderef{TashkinovSeriesDataAppend}\).  For every old
non-tree index \(i\),
\[
  D_{A^+}(i^+)=D_A(i),
\]
where \(i^+\) is the canonical cast-successor embedding and both sides use
the stored-color definition of \(\noderef{TashkinovStoredColors}\).
\end{lemma}
\begin{proof}
The cast-successor index \(i^+\) has the same numerical value as \(i\).  If
that value is zero, both sides are empty by the explicit zero clause in
\(\noderef{TashkinovStoredColors}\).

Suppose it is positive.  The two set differences remove the same tree-missing
colors: by \(\noderef{TashkinovSeriesDataAppend}\), the coloring at \(i^+\)
is the old coloring at \(i\), and the tree at the cast-successor of \(i^+\)
is the old tree at the cast-successor of \(i\).

It remains to compare the bounded unions of stored sets.  If an enlarged
index \(h\in\operatorname{Fin}(r+2)\) satisfies \(h\leq i^+\), then
\(h.val\leq i.val<i.val+1\leq r+1\).  Hence
\(h_0=\langle h.val,h.val<r+1\rangle\) is an old index and
\(h=h_0^+\).  The append definition identifies its stored set with the old
stored set at \(h_0\).  This converts every member of the enlarged bounded
union into a member of the old bounded union.  Conversely, an old index
\(h\leq i\) maps to \(h^+\leq i^+\), and its stored set is unchanged, giving
the reverse inclusion.  Thus the bounded unions, the removed missing-color
sets, and finally the two differences are equal.
\end{proof}
